\section{总结与延伸}

\subsection{核心知识图谱}

本节课建立了从词表示到大语言模型的完整认知链：

\begin{center}
\begin{tikzpicture}[node distance=0.9cm, every node/.style={draw, rounded corners, minimum width=4cm, minimum height=0.7cm, font=\small}]
\node (sub) {子词分词 (BPE)};
\node (ctx) [below=of sub] {上下文表示 vs 静态嵌入};
\node (para) [below=of ctx] {预训练--微调范式};
\node (enc) [below=of para] {Encoder: BERT (MLM)};
\node (encdec) [below=of enc] {Enc-Dec: T5 (Span Corruption)};
\node (dec) [below=of encdec] {Decoder: GPT (LM)};
\node (scale) [below=of dec] {规模化与涌现能力};
\draw[->, thick] (sub) -- (ctx);
\draw[->, thick] (ctx) -- (para);
\draw[->, thick] (para) -- (enc);
\draw[->, thick] (enc) -- (encdec);
\draw[->, thick] (encdec) -- (dec);
\draw[->, thick] (dec) -- (scale);
\node[draw=none, right=0.8cm of sub, text width=5cm, font=\scriptsize] {解决 OOV 问题，词表大小 vs 序列长度的权衡};
\node[draw=none, right=0.8cm of ctx, text width=5cm, font=\scriptsize] {Word2Vec 的一词一向量 vs Transformer 的上下文向量};
\node[draw=none, right=0.8cm of para, text width=5cm, font=\scriptsize] {海量无标注数据 $\rightarrow$ 少量标注数据微调};
\node[draw=none, right=0.8cm of enc, text width=5cm, font=\scriptsize] {双向上下文，擅长理解，不能生成};
\node[draw=none, right=0.8cm of encdec, text width=5cm, font=\scriptsize] {理解+生成，序列到序列任务};
\node[draw=none, right=0.8cm of dec, text width=5cm, font=\scriptsize] {单向上下文，擅长生成，规模化的主流选择};
\node[draw=none, right=0.8cm of scale, text width=5cm, font=\scriptsize] {ICL, CoT, Scaling Laws, Chinchilla};
\end{tikzpicture}
\end{center}

\subsection{关键 Takeaways}

\begin{importantbox}{本课七条核心原则}
\begin{enumerate}
    \item \textbf{子词分词解决了开放词表问题}：BPE 等算法使模型能够处理任意输入，消除了 UNK 标记
    \item \textbf{上下文表示优于静态嵌入}：同一个词在不同上下文中应有不同的向量表示
    \item \textbf{预训练--微调是当今 NLP 的基础范式}：先在海量无标注文本上学习通用表示，再用少量标注数据适配具体任务
    \item \textbf{三种架构各有所长}：Encoder（理解）、Encoder-Decoder（理解+生成）、Decoder（生成），但 Decoder-only 正成为主流
    \item \textbf{规模带来涌现能力}：上下文学习、Chain-of-Thought 推理等能力在大规模模型中``意外''出现
    \item \textbf{计算效率很重要}：Chinchilla 表明参数量与数据量应等比例缩放，盲目增大模型不一定最优
    \item \textbf{能力与风险并存}：预训练模型学到的知识包括有害偏见，部署时必须谨慎评估
\end{enumerate}
\end{importantbox}

\subsection{拓展阅读}

\begin{itemize}
    \item Devlin et al., \textbf{BERT: Pre-training of Deep Bidirectional Transformers for Language Understanding} (2018): \url{https://arxiv.org/abs/1810.04805}
    \item Liu et al., \textbf{RoBERTa: A Robustly Optimized BERT Pretraining Approach} (2019): \url{https://arxiv.org/abs/1907.11692}
    \item Joshi et al., \textbf{SpanBERT: Improving Pre-training by Representing and Predicting Spans} (2020): \url{https://arxiv.org/abs/1907.10529}
    \item Raffel et al., \textbf{Exploring the Limits of Transfer Learning with a Unified Text-to-Text Transformer} (T5, 2020): \url{https://arxiv.org/abs/1910.10683}
    \item Radford et al., \textbf{Improving Language Understanding by Generative Pre-Training} (GPT, 2018)
    \item Brown et al., \textbf{Language Models are Few-Shot Learners} (GPT-3, 2020): \url{https://arxiv.org/abs/2005.14165}
    \item Hoffmann et al., \textbf{Training Compute-Optimal Large Language Models} (Chinchilla, 2022): \url{https://arxiv.org/abs/2203.15556}
    \item Wei et al., \textbf{Chain-of-Thought Prompting Elicits Reasoning in Large Language Models} (2022): \url{https://arxiv.org/abs/2201.11903}
    \item Hu et al., \textbf{LoRA: Low-Rank Adaptation of Large Language Models} (2021): \url{https://arxiv.org/abs/2106.09685}
    \item Gururangan et al., \textbf{Don't Stop Pretraining} (2020): \url{https://arxiv.org/abs/2004.10964}
\end{itemize}

\end{document}
